\centerline{\bf \Huge Оценка + Пример I}

\begin{flushright}
        \scriptsize{Твои возможности обычно больше, чем твоя их оценка\\ Шри Ауробиндо}
\end{flushright}
\problem{\textbf{1}} В гости пришло 10 гостей и каждый оставил в коридоре пару калош. Все пары калош имеют разные размеры. Гости начали расходиться по одному, одевая любую пару калош, в которые они могли влезть (т.е. каждый гость мог надеть пару калош, не меньшую, чем его собственные). В какой-то момент обнаружилось, что ни один из оставшихся гостей не может найти себе пару калош, чтобы уйти. Какое максимальное число гостей могло остаться?

\problem{\textbf{2}} В сумме  + 1 + 3 + 9 + 27 + 81 + 243 + 729  можно вычеркивать любые слагаемые и изменять некоторые знаки перед оставшимися числами с "+" на "–". Маша хочет таким способом сначала получить выражение, значение которого равно 1, затем, начав сначала, получить выражение, значение которого равно 2, затем (снова начав сначала) получить 3, и так далее. До какого наибольшего целого числа ей удастся это сделать без пропусков?


\problem{\textbf{3}} Наиль расставляет в клетках квадрата 6×6 числа от 1 до 36 (по одному числу в каждую клетку, числа не повторяются). После этого Наиль ставит фишку в клетку с числом 1. Далее перед каждым ходом Наиль выбирает наибольшее из чисел, стоящих в соседних с фишкой (по стороне или углу) клетках. Если выбранное число больше, чем в клетке с фишкой, то Наиль передвигает фишку в клетку с выбранным числом; иначе фишка больше не двигается. При какой изначальной расстановке чисел фишка посетит как можно больше клеток?

\problem{\textbf{4}} В каждой клетке таблицы $5\times5$ записано число. Назовём клетку $C$ хорошей, если в какой-то из клеток, соседних с $C$ по стороне, стоит число на 1 больше, чем в $C$, а в какой-то другой из клеток, соседних с $C$ по стороне, стоит число на 3 больше, чем в $C$. Каково наибольшее возможное количество хороших клеток?


\problem{\textbf{5}} В некоторых клетках доски  $10\times10$ поставили $k$  ладей, и затем отметили все клетки, которые бьет хотя бы одна ладья (считается, что ладья бьет клетку, на которой стоит). При каком наибольшем  $k$ может оказаться, что после удаления с доски любой ладьи хотя бы одна отмеченная клетка окажется не под боем?


\centerline{\textit{1 балл за Пример + 1 балл за Оценку}}\\

